% main.tex -- INDIACom-2027 (Track 4: Software Engineering and Emerging Technologies)
% DRAFT FOR AUTHOR AND GUIDE REVIEW. Every number is traced in EVIDENCE_LEDGER.md.
% Passages marked [AUTHOR INPUT] must be rewritten in the authors' own words.
\documentclass[conference,a4paper]{IEEEtran}
\IEEEoverridecommandlockouts
\usepackage{cite}
\usepackage{amsmath}
\usepackage{graphicx}
\usepackage{booktabs}
\usepackage{multirow}
\usepackage{url}
\usepackage{xcolor}
\usepackage{tikz}
\usetikzlibrary{positioning,arrows.meta,fit}

% Set \anonymoustrue to build a double-blind review copy (see CONFERENCE_REQUIREMENTS.md).
\newif\ifanonymous
\anonymousfalse

\newcommand{\killk}[1]{Kill@#1}
\newcommand{\authorinput}[1]{\textcolor{red}{[AUTHOR INPUT: #1]}}

\begin{document}

\title{Oneiros Engine: A Self-Learning Framework for Intelligent and Realistic Software Testing}

\ifanonymous
\author{\IEEEauthorblockN{Anonymous Authors}\IEEEauthorblockA{Affiliation withheld for review}}
\else
\author{
\IEEEauthorblockN{[AUTHOR NAME 1]}
\IEEEauthorblockA{[USN]\\
Department of Computer Science and Engineering\\{}
[AFFILIATION -- TO VERIFY: PES University]\\{}
[CITY -- TO VERIFY: Bengaluru], India\\{}
[EMAIL]}
\and
\IEEEauthorblockN{[AUTHOR NAME 2]}
\IEEEauthorblockA{[USN]\\
Department of Computer Science and Engineering\\{}
[AFFILIATION -- TO VERIFY: PES University]\\{}
[CITY -- TO VERIFY: Bengaluru], India\\{}
[EMAIL]}
\and
\IEEEauthorblockN{[GUIDE NAME]}
\IEEEauthorblockA{Department of Computer Science and Engineering\\{}
[AFFILIATION -- TO VERIFY: PES University]\\{}
[CITY -- TO VERIFY: Bengaluru], India\\{}
[EMAIL]}
}
\fi

\maketitle

\begin{abstract}
A generated unit test is useful only if it chooses an input that exposes a fault
and also states the correct expected output. Mutation-killing rates reward the first
property but can hide failures of the second. We present Oneiros, a framework that
fine-tunes a small code language model to write mutation-killing tests and evaluates
every candidate by running the same test on the fixed and on the faulty program, under
a frozen, hash-bound protocol. Completion-only QLoRA fine-tuning of
Qwen2.5-Coder-1.5B raised the share of development targets killed by at least one of
eight candidates from 60.52 to 69.56 percent. On a locked validation split never used
for selection, the gain shrank to 3.43 points (59.45 to 62.88 percent; exact McNemar
p-value 0.078, above the predeclared 0.01 threshold), while the share of
reference-valid candidates fell from 45.21 to 33.34 percent and wrong-oracle candidates
rose from 1,863 to 2,736 of 6,056. The checkpoint was therefore not promoted. The
evidence indicates that fine-tuning improved input selection and assertion form, not
expected-value prediction, and that kill rate alone is an unsafe promotion metric.
\end{abstract}

\begin{IEEEkeywords}
Software testing, large language models, mutation testing, supervised fine-tuning,
test generation, oracle correctness, empirical software engineering
\end{IEEEkeywords}

%=====================================================================================
\section{Introduction}

A unit test has two jobs. It must drive the program to a state where a fault shows,
and it must state what the correct program should produce there. The second job is the
test-oracle problem~\cite{barr2015oracle}. Large language models (LLMs) now generate
plausible unit tests~\cite{schafer2024llmunit,lemieux2023codamosa,wang2024testingsurvey},
and mutation testing~\cite{jia2011mutation,papadakis2019mutation} is a common way to
judge them: a test that fails on a seeded mutant but passes on the original
``kills'' it.

A kill rate is easy to optimise. A test that asserts \emph{some} value on a
discriminating input kills the mutant whenever the mutant's value differs from the
asserted one, even if the asserted value is also wrong for the correct program. If a
model is selected only on kill rate, it can learn to choose good inputs while getting
worse at predicting outputs. That regression is invisible unless every candidate is
also run on the correct program.

This paper reports what happened when we fine-tuned a 1.5B-parameter code model for
mutation-killing test generation and measured both properties under a frozen
protocol. We frame it as an empirical diagnostic study, not a performance claim. We ask:

\begin{itemize}
\item \textbf{RQ1:} Does completion-only QLoRA fine-tuning improve mutation killing
under a frozen \killk{8} protocol?
\item \textbf{RQ2:} Do \killk{8} gains preserve reference validity and correct
oracle prediction?
\item \textbf{RQ3:} Which failure categories explain the gap between development
gains and locked-validation behaviour?
\item \textbf{RQ4 (exploratory):} To what extent can these conclusions currently be
checked on repository-native functions?
\end{itemize}

\textbf{Contributions.} (1)~A source-bound, leakage-controlled pipeline for training
and evaluating LLM-generated mutation-killing tests, with hash-bound receipts for data,
model, protocol and outputs. (2)~Whole-output candidate retention and dual execution
of each candidate on fixed and buggy code. (3)~A controlled base-versus-fine-tuned
comparison on development and locked-validation splits. (4)~A failure taxonomy that
separates valid kills from wrong oracles, non-discriminating inputs, parse and
execution failures, and duplicates. (5)~Evidence that \killk{8} alone is not a safe
promotion metric, and a multi-gate promotion procedure. (6)~A repository-native
evaluation pipeline, reported here as built but not yet executed.

\authorinput{add one or two sentences on why the team started Oneiros and what
practical testing problem motivated it.}

%=====================================================================================
\section{Related Work}

\textbf{Automated and LLM-based test generation.} Search-based and random tools such
as EvoSuite~\cite{fraser2011evosuite} and Randoop~\cite{pacheco2007randoop} generate
tests and use regression oracles taken from the current program. CodaMosa uses an LLM
to escape coverage plateaus in search-based generation~\cite{lemieux2023codamosa}, and
Sch\"afer \emph{et al.} evaluate LLM unit-test generation directly~\cite{schafer2024llmunit};
Wang \emph{et al.} survey the area~\cite{wang2024testingsurvey}. MuTAP feeds surviving
mutants back into LLM prompts to improve mutation score~\cite{dakhel2024mutap}.

\textbf{Mutation testing.} Mutation analysis is a standard adequacy
criterion~\cite{jia2011mutation,papadakis2019mutation}, and mutant detection
correlates with real-fault detection~\cite{just2014mutants}. A kill, however, is
evidence of a correct oracle only when the test also passes on the correct program.

\textbf{Oracles.} Barr \emph{et al.} survey the oracle problem~\cite{barr2015oracle}.
TOGA trains a neural model to generate assertion oracles~\cite{dinella2022toga}. Our
study measures oracle correctness as a separate outcome rather than assuming it.

\textbf{Fine-tuning and benchmarks.} We adapt Qwen2.5-Coder~\cite{hui2024qwen25coder}
with LoRA~\cite{hu2021lora} on a 4-bit quantised base (QLoRA)~\cite{dettmers2023qlora}.
Our function-level targets derive from HumanEval~\cite{chen2021codex} and
MBPP~\cite{austin2021mbpp}; EvalPlus showed that weak test suites overstate code
correctness~\cite{liu2023evalplus}. Repository targets derive from SWE-bench
Verified~\cite{jimenez2023swebench} and BugsInPy~\cite{widyasari2020bugsinpy}.

\textbf{Fuzzing and evaluation validity.} Coverage-guided fuzzers such as
libFuzzer~\cite{serebryany2016libfuzzer} and its Python front end
Atheris~\cite{atheris} are natural baselines, but only when correctly instrumented.
Data leakage is a known source of inflated ML results~\cite{kapoor2023leakage}, which
motivates our lineage-grouped splits and single-use locked evaluation.

%=====================================================================================
\section{Oneiros Architecture and Methodology}

Fig.~\ref{fig:workflow} shows the workflow.

\begin{figure*}[t]
\centering
\input{figures/oneiros_workflow.tex}
\caption{Oneiros workflow. Data and training stages (top) produce a checkpoint chosen
only on the development split; evaluation stages (bottom) execute each candidate on
both the fixed and the faulty program, classify it, and apply a multi-gate promotion
rule. Every stage writes a hash-bound receipt.}
\label{fig:workflow}
\end{figure*}

\textbf{Verified program pairs.} Each record pairs a correct (fixed) implementation
with a faulty one. Function-level records are mutants of benchmark solutions, grouped
by mutation family (e.g.\ boundary, arithmetic, comparison). Repository records are
real defects with official reproductions from SWE-bench Verified and BugsInPy. A record
is admitted only if at least one reference-valid, mutation-killing test exists for it
under the current policy.

\textbf{Normalisation and lineage-controlled splitting.} Records are normalised to one
prompt schema and split by semantic group, so all mutants of one original function
fall in the same split. A leakage audit checks for verbatim reference code in prompts
and for lineage links from gold tests, gold patches and oracles.

\textbf{Prompt and target.} The prompt contains the task specification and the code
under test, within fixed token budgets. The supervised target is a verified
mutation-killing test that has already been executed against both versions.

\textbf{Completion-only QLoRA fine-tuning.} The base model is loaded in 4-bit NF4 and
frozen; low-rank adapters are trained on the completion tokens only. Checkpoints are
selected on the development split; the locked validation split is read once, after
selection, under a decision rule frozen beforehand.

\textbf{Generation and parsing.} For each target the model emits eight ordered
candidates with no reranking. Each raw output is retained and hashed, and the whole
output is parsed as one test (the \emph{whole-output} protocol; see
Section~\ref{sec:threats}). A safety policy rejects unsafe constructs before execution.

\textbf{Dual execution and classification.} The \emph{same} candidate runs in a
sandbox on the fixed and on the faulty implementation. Each candidate receives exactly
one terminal category, by fixed precedence: \emph{parse failure};
\emph{duplicate} (whitespace-normalised code equal to an earlier candidate for the same
target); \emph{valid kill} (passes on fixed, fails on faulty); \emph{execution
failure} (does not reach its assertions on the fixed code); \emph{non-discriminating
input} (no call behaves differently on the mutant); \emph{discriminating wrong oracle}
(a call behaves differently, but an assertion fails on the correct code); and
\emph{valid non-kill} (correct but too weak to fail on the mutant). A target is
killed at $k$ if any of its first $k$ candidates is a valid kill (\killk{k}).

\textbf{Receipts.} Corpus, splits, model revision, adapter, protocol, source tree,
raw outputs, metrics and decisions are bound by SHA-256 digests. A full-scale
rehearsal of the evaluation path reproduced the recorded base-model development
\killk{8} exactly (0.605166) over 4,336 candidates with zero missing or mismatched
raw-output hashes, which validates the pipeline, not the model.

%=====================================================================================
\section{Dataset and Experimental Setup}

\begin{table}[t]
\centering
\caption{Corpus V4.1: Sources, Split Roles and Complexity}
\label{tab:data}
\footnotesize
\setlength{\tabcolsep}{3pt}
\begin{tabular}{@{}lrl@{}}
\toprule
\multicolumn{3}{@{}l}{\textit{Records by source (8,237 total; 7,644 function, 593 repository)}}\\
MBPP & 6,384 & synthetic mutants\\
HumanEval & 1,252 & synthetic mutants\\
SWE-bench Verified & 324 & real repository defects\\
BugsInPy & 269 & real repository defects\\
Manually curated & 8 & seeded examples\\
\midrule
\multicolumn{3}{@{}l}{\textit{Splits (records / semantic groups) and role}}\\
Train & 6,052 / 683 & training pool only\\
Development & 594 / 63 & checkpoint selection (542 evaluated)\\
Locked validation & 781 / 82 & read once (757 evaluated)\\
Sealed test & -- & consumed by a failed run; no metrics\\
\midrule
\multicolumn{3}{@{}l}{\textit{Fine-tuning selection (2,400 pairs)}}\\
Sources & \multicolumn{2}{l}{MBPP 1,776; HumanEval 413; SWE-bench 132;}\\
 & \multicolumn{2}{l}{BugsInPy 74; manual 5}\\
Complexity (synthetic) & \multicolumn{2}{l}{complex 1,317; moderate 610; simple 267}\\
Effective examples & \multicolumn{2}{l}{6,852 (repository share 0.20)}\\
\midrule
\multicolumn{3}{@{}l}{\textit{Repository-native panel (exploratory; cx.\,=\,complex)}}\\
Frozen panel & 34 & 25 cx., 7 mod., 2 simple; 15 repos\\
Native-executable & 29 & 22 cx., 5 mod., 2 simple; 14 repos\\
\bottomrule
\end{tabular}
\end{table}

\textbf{Corpus and roles.} Table~\ref{tab:data} summarises corpus V4.1. Repository
records are excluded from every kill rate reported in Section~\ref{sec:results},
because they need a native project environment; only function-mode targets are scored
(542 on development, 757 on locked validation). The leakage audit scanned all 8,237
records: zero verbatim reference leaks, zero gold-test, gold-patch or oracle lineage
failures, and 14 partial-overlap flags, all manually reviewed.

\textbf{Complexity.} Function tiers come from the AST of the faulty function:
\emph{complex} if any of cyclomatic complexity $\geq 6$, control nesting $\geq 3$,
logical lines $\geq 24$ or AST nodes $\geq 80$ holds; \emph{moderate} at
$3/2/12/40$; otherwise \emph{simple}. Repository regions use larger thresholds plus a
branch count. Complex functions make up 1,317 of the 2,194 synthetic training pairs,
and 25 of the 34 repository-native targets.

\textbf{Model and training.} The base model is Qwen2.5-Coder-1.5B-Instruct at
immutable revision \texttt{2e1fd397}. The adapter (``A@431'', the checkpoint at
optimiser step 431) uses LoRA rank 16, $\alpha=32$, dropout 0.05 on the query, key,
value and output projections, with NF4 4-bit loading, learning rate $10^{-5}$
(constant after 25 warm-up steps), one epoch, and 16 examples per optimiser step.
Selected pairs are re-weighted so that repository examples form 20\% of the 6,852
effective examples.

\textbf{Frozen evaluation protocol.} Eight candidates per target, temperature 0.7,
top-$p$ 0.9, seed 42, whole-output parsing, 1,024 completion and 3,072 sequence
tokens, raw outputs retained, no reranking. Intervals for single rates are Wilson
95\%~\cite{wilson1927probable}; paired comparisons use the exact McNemar
test~\cite{mcnemar1947note} over targets and a paired bootstrap (10,000 resamples)
clustered by semantic group.

\textbf{Predeclared promotion rule.} Promotion required all of: (1)~a locked
\killk{8} gain of at least $+3.0$ points, a net gain of at least 15 targets, and
McNemar $p<0.01$; (2)~a reference-validity drop of at most 12 points; (3)~parse and
execution drops of at most 1 point each; (4)~no diversity collapse; and (5)~run
integrity.

%=====================================================================================
\section{Results}
\label{sec:results}

\begin{table}[t]
\centering
\caption{Base vs. Fine-Tuned (A@431). Rates in Percent}
\label{tab:main}
\footnotesize
\setlength{\tabcolsep}{2.6pt}
\begin{tabular}{@{}lrrrr@{}}
\toprule
 & \multicolumn{2}{c}{Development ($n$=542)} & \multicolumn{2}{c}{Locked val. ($n$=757)}\\
\cmidrule(lr){2-3}\cmidrule(l){4-5}
Metric & Base & A@431 & Base & A@431\\
\midrule
\killk{1} & 23.43 & 34.13 & 23.65 & 23.91\\
\killk{4} & 48.34 & 57.93 & 47.29 & 47.95\\
\killk{8} & \textbf{60.52} & \textbf{69.56} & \textbf{59.45} & \textbf{62.88}\\
Targets killed @8 & 328 & 377 & 450 & 476\\
\killk{8} Wilson CI & [56.3, 64.5] & [65.6, 73.3] & [55.9, 62.9] & [59.4, 66.2]\\
Parse success & 98.34 & 99.82 & 98.73 & 98.88\\
Execution-valid & 94.86 & 99.31 & 92.83 & 94.22\\
Reference-valid & \textbf{56.55} & \textbf{45.50} & \textbf{45.21} & \textbf{33.34}\\
\midrule
$\Delta$\killk{8} (points) & \multicolumn{2}{c}{$+9.04$} & \multicolumn{2}{c}{$+3.43$}\\
Gained / lost targets & \multicolumn{2}{c}{100 / 51} & \multicolumn{2}{c}{114 / 88}\\
Exact McNemar $p$ & \multicolumn{2}{c}{$8.2\times10^{-5}$\,$^\dagger$} & \multicolumn{2}{c}{0.0783}\\
Clustered CI, $\Delta$\killk{8} & \multicolumn{2}{c}{--} & \multicolumn{2}{c}{[$-4.03$, $+10.88$]}\\
$\Delta$ ref.-valid (points) & \multicolumn{2}{c}{$-11.05$} & \multicolumn{2}{c}{$-11.87$}\\
\bottomrule
\multicolumn{5}{@{}p{0.95\columnwidth}@{}}{\scriptsize $^\dagger$The development split
selected the checkpoint, so this value is optimistically biased and is not
generalisation evidence.}
\end{tabular}
\end{table}

\textbf{RQ1 -- mutation killing.} Table~\ref{tab:main} gives the results. On the
development split, A@431 raised \killk{8} from 60.52\% to 69.56\% ($+9.04$ points;
100 targets gained, 51 lost). Because that split chose the checkpoint, this is a
selection statistic. On locked validation, read once, the gain was $+3.43$ points
(59.45\% to 62.88\%; 114 gained, 88 lost, net $+26$). The exact McNemar $p$ was 0.0783,
and the 95\% interval clustered over 78 semantic groups, $[-4.03, +10.88]$, includes
zero. The minimum effect detectable with 80\% power at this cluster count is about
10.7 points. \killk{1} and \killk{4} barely moved on locked validation ($+0.26$ and
$+0.66$ points). The development gain therefore transferred only partly, and the
locked gain is not statistically established.

\textbf{RQ2 -- validity.} Reference validity, the share of candidates that pass on the
correct program, fell by 11.05 points on development and by 11.87 points on locked
validation (45.21\% to 33.34\%; clustered CI $[-20.09, -3.26]$). Targets with at least
one reference-valid candidate fell from 675 to 566 of 757. Parse and execution rates
did not fall, so the regression is in test \emph{content}, not syntax or runnability.

\begin{table}[t]
\centering
\caption{Terminal Categories on Locked Validation (6,056 Candidates per Arm)}
\label{tab:taxonomy}
\footnotesize
\setlength{\tabcolsep}{4pt}
\begin{tabular}{@{}lrrrr@{}}
\toprule
 & \multicolumn{2}{c}{Base} & \multicolumn{2}{c}{A@431}\\
\cmidrule(lr){2-3}\cmidrule(l){4-5}
Category & Count & \% & Count & \%\\
\midrule
Valid kill & 975 & 16.10 & 1,105 & 18.25\\
Discriminating wrong oracle & 1,863 & 30.76 & \textbf{2,736} & \textbf{45.18}\\
Non-discriminating input & 1,226 & 20.24 & 1,127 & 18.61\\
Valid non-kill (weak oracle) & 497 & 8.21 & 1 & 0.02\\
Duplicate & 1,118 & 18.46 & 817 & 13.49\\
Execution failure & 300 & 4.95 & 202 & 3.34\\
Parse failure & 77 & 1.27 & 68 & 1.12\\
\bottomrule
\end{tabular}
\end{table}

\begin{table}[t]
\centering
\caption{Stage Rates on Locked Validation (Percent; $\Delta$ with 95\% CI Clustered over 78 Semantic Groups)}
\label{tab:stages}
\footnotesize
\setlength{\tabcolsep}{2.5pt}
\begin{tabular}{@{}lrrr@{}}
\toprule
Stage & Base & A@431 & $\Delta$ [95\% CI]\\
\midrule
Discriminating $\mid$ executed & 73.37 & 77.74 & $+4.37$ [1.94, 6.94]\\
Correct oracle $\mid$ discriminating & 47.03 & 33.18 & $-13.85$ [$-22.52$, $-4.80$]\\
\quad same, equality asserts only$^\ast$ & 37.39 & 33.14 & $-4.25$ [$-11.09$, 2.67]\\
Equality-form assertions & 72.47 & 99.50 & $+27.03$ [18.56, 35.62]\\
Kill $\mid$ reference-valid & 50.95 & 72.71 & $+21.76$ [11.73, 31.82]\\
\bottomrule
\multicolumn{4}{@{}p{0.95\columnwidth}@{}}{\scriptsize $^\ast$Form-stratified metric
added after the first analysis; reported alongside, not instead of, the predeclared
unstratified row.}
\end{tabular}
\end{table}

\textbf{RQ3 -- failure categories.} Table~\ref{tab:taxonomy} shows where the
candidates went. Fine-tuning added 130 valid kills but 873 discriminating wrong
oracles, and nearly eliminated weak (valid non-kill) tests. Table~\ref{tab:stages}
decomposes the pipeline. The fine-tuned model chose discriminating inputs more often
($+4.37$ points) and, once its test was valid, killed far more often ($+21.76$). But
given a discriminating input, its assertion held on the correct program much less
often ($-13.85$). The model also moved almost all assertions to exact equality
(72.47\% to 99.50\%). Restricted to equality assertions, the change in oracle
correctness was $-4.25$ points with an interval that includes zero. The more cautious
reading is that fine-tuning transferred \emph{input choice and assertion form} but not
the ability to predict the correct value: sharper assertions turn uncertain value
predictions into both more kills and more wrong oracles. At the target level, 69
targets gained and 150 lost a valid discriminating candidate.

\textbf{Promotion decision.} Criterion~1 failed on significance ($p=0.0783$ against
$p<0.01$), although its gain ($\geq +3.0$) and net ($\geq 15$) parts passed. Criterion~2
\emph{passed}: the 11.87-point validity drop was inside the predeclared 12-point
allowance. Criterion~5 failed because of an implementation defect: a late change put
the adapter hash inside the run-contract digest, so a fine-tuned arm could never match
the base arm. This was recorded as a failure, not reinterpreted. The decision was to
retain the base model.

\textbf{RQ4 -- repository-native evidence.} From 70 frozen repositories, 52 were
scanned and 804 candidate defects evaluated; 34 targets survived native qualification,
requalification, reach checks and prompt-leakage and token-fit checks. Of these, 29
from 14 repositories are executable under the sandbox policy; 5 were excluded before
any generation because the policy blocks native-library loading. The predeclared
minimum for a confirmatory panel was 80 targets from 10 repositories, so the panel is
recorded as underpowered and exploratory: at 29 targets, a planning estimate puts the
minimum detectable effect near 26 points. A three-arm generation plan (base, A@431 and
a relearning checkpoint) with a frozen analysis passed all 17 preflight checks, but
\textbf{no repository-native generation has yet been run}. We report no
repository-native result.

%=====================================================================================
\section{Discussion}

\textbf{Input selection and oracle prediction are separate skills.} A test must both
reach a state where the fault is visible and predict the correct output there. Our
measurements separate the two. Fine-tuning helped the first and did not help the
second. Because a kill requires only that the faulty program disagrees with the
asserted value, a model that asserts \emph{some} value more forcefully on a better
input can raise \killk{8} while more of its tests are wrong for the correct program.
\killk{8} can rise while semantic validity falls.

\textbf{Why the checkpoint was not promoted.} The locked gain was positive but within
noise, and the development gain was optimistic because of selection. A rule based on
\killk{8} alone, especially on the selection split, would have promoted a model that
produces 873 more wrong-oracle candidates on unseen targets.

\textbf{Lessons for promotion rules.} Our frozen rule allowed a validity drop of up to
12 points, and the observed drop (11.87) passed. In hindsight that allowance was too
loose: a promotion gate should treat validity and wrong-oracle growth as
\emph{blocking}, not as tolerances around a kill-rate gain.

\textbf{Implications.} Studies that report mutation score or pass rate for
LLM-generated tests should also report how many candidates fail on the correct
program, and should select checkpoints on one split and report on another.
\authorinput{add the team's own interpretation and any implementation experience
that informed these lessons.}

%=====================================================================================
\section{Threats to Validity and Limitations}
\label{sec:threats}

\textbf{Scale and data.} One 1.5B model and one training seed were evaluated. Training
data are dominated by MBPP and HumanEval mutants; real repository defects form a
small share. Results may not hold for larger models, other languages or other fault
types.

\textbf{Statistical power.} With 757 targets in 78 semantic groups, the locked study
could reliably detect only gains of about 10.7 points; a non-significant result is not
evidence of no effect. The locked split has now been used and cannot be re-measured.

\textbf{Sealed test.} The single authorised sealed-test run failed during record
admission, before any generation. It produced no candidates and no metrics, and the
split is consumed. No final-test result exists for Oneiros.

\textbf{Evaluator defects found and corrected.} An earlier parser scored only the first
\texttt{assert} line of each output; the whole-output protocol used here replaced it,
and all results in this paper use the corrected protocol. An earlier Atheris baseline
was not coverage-instrumented and passed aliased arguments to reference and mutant; it
was corrected, but its figures were measured under the earlier test protocol, so we do
not compare Oneiros with Atheris here. The integrity-criterion defect is described in
Section~\ref{sec:results}.

\textbf{Measurement choices.} One terminal category per candidate depends on a fixed
precedence order. The form-stratified oracle metric was added after the first look.
Some HumanEval prompts state an input--output example that kills the mutant, which
inflates HumanEval kill rates. All execution depends on the sandbox policy, which also
excluded 5 repository targets.

\textbf{Repository-native evaluation} remains underpowered and unexecuted.

%=====================================================================================
\section{Conclusion and Future Work}

Fine-tuning a small code model on verified mutation-killing tests raised development
\killk{8} by 9.04 points, but on a locked split the gain was 3.43 points and not
significant, while reference validity fell by 11.87 points. The failure taxonomy
locates the cause: better input selection and sharper assertions, without better
expected-value prediction. Oneiros's main contribution is therefore the evaluation and
promotion methodology that caught this, not a better model.

Next steps are to: treat input selection and oracle prediction as separate learning
targets; make validity and wrong-oracle growth blocking promotion gates; build more
unique, execution-verified repository examples and keep replay data during
fine-tuning; complete a sufficiently powered repository-native evaluation on a new,
independently built held-out set; and compare against corrected baselines on identical
eligible subsets.

\ifanonymous\else
\section*{Acknowledgment}
{\raggedright\authorinput{acknowledge the guide and the department; review the
AI-assistance statement below and keep it accurate.}\par}
An AI assistant (Anthropic Claude, used through Claude Code) helped organise this
manuscript and produced an initial draft of its text, tables and Fig.~1 from the
project's committed experiment records. All authors reviewed, revised and verified
every statement, number and citation, and take full responsibility for the content.
\fi

\bibliographystyle{IEEEtran}
\bibliography{references}

\end{document}
